\section{Experiment}
\subsection{Experimental Setup}
\paragraph{Models.}
We use the publicly released Wan2.2-TI2V-5B dense video diffusion model in its
T2V mode. Videos contain 121 frames at $1280\times704$ resolution and 24 FPS.
All methods use the same model, 50-step UniPC sampling configuration, random-seed
protocol, and decoding pipeline; mixed-precision inference and untimed CUDA
warm-up follow the formal runner.

\paragraph{Hardware.}
Experiments run on a server equipped with eight NVIDIA H100 GPUs. Parallel
baseline workers are assigned one GPU each; the remaining devices are available
for cache preparation, dense-control runs, and quality evaluation.

\paragraph{Datasets.}
We use VBench and VidProM. VBench contributes 694 deduplicated prompts from eight
dimensions; we report CLIP-T and the six custom-input dimensions used for
VBench-Q. VidProM prompts are filtered for usability, after which 1,000 requests
are randomly selected as preloaded cache entries and 300 non-overlapping requests
as new arrivals. For object-level reuse, each target has a cache candidate under
the noun-match relation (identical noun sets or one noun substitution); targets
follow a fixed arrival order and are not inserted into the source cache.

\paragraph{Metrics.}
For quality, we report CLIP-T, computed from eight uniformly sampled frames, and
VBench-Q over the six selected dimensions. For efficiency, we measure end-to-end
and CUDA-synchronized DiT latency, executed steps, hit rate, hit latency, peak
GPU memory, and cache size. Chorus's prompt-difference and segmentation overhead
is reported separately but included in end-to-end latency. Speedups use the
same-configuration dense baseline on the same 300 requests.

\paragraph{Baselines and variants.}
The formal VidProM experiment compares one dense baseline with five acceleration
methods. Each method uses a persistent single-GPU worker and processes the same
300 requests in the same arrival order.

\begin{itemize}
  \item \textbf{Dense baseline.} Generates every request with the complete
  50-step dense denoising process and provides the quality and speed reference.
  \item \textbf{NIRVANA-VID.} Retrieves a semantically similar cached request
  using full-prompt embeddings and reuses its intermediate latent state to
  restart the target denoising process from an early step.
  \item \textbf{Chorus.} Retrieves a similar request, reuses its denoising
  trajectory, and applies regional updates with TGAA. In our reproduction, the
  prompt difference is extracted by Qwen3.5-Plus (qwen3.5-plus-2026-02-15) through an OpenAI-compatible API. Object masks are produced by Grounded-SAM-2 v1.0 using Grounding DINO Swin-T OGC and SAM 2.1 Hiera-Large. The LLM is instructed to compare source and target prompts and
  return JSON fields for source entities, target entities, removed tokens, added tokens, and an edit summary. Grounded-SAM-2 is queried with the extracted source-entity phrase (up to four unique entities). The associated LLM and segmentation times are reported as external overhead.
  \item \textbf{Chorus II.} Reuses cached spatial attention information from a
  similar request while continuing target denoising, without restarting from a
  cached latent state. The implementation uses the source block masks with
  sparse attention in the reuse stage.
  \item \textbf{Carnator.} First narrows candidates using noun-level matching,
  then applies an Early Signature compatibility test and embedding retrieval.
  For accepted pairs, it builds a target-specific regional mask from cached
  per-noun artifacts and accelerates the boundary denoising stage with sparse
  attention and W8A8 FFN computation. Failed matching or mask validation is
  recorded as a miss.
\end{itemize}

Beyond these method-specific settings, all experiments share the same source-cache
boundary, input order, model sampling parameters, and quality-evaluation
pipeline. The source cache is prepared once and exposed read-only to every
worker; each method maintains independent retrieval state and an independent
result directory.




\subsection{Main Results}

The current run has completed all five acceleration workers on the 300-request
VidProM workload, but the matched Dense C0 control is still running.  Table~\ref{tab:main_results}
therefore reports the measured values from the completed run and marks the dense
reference and the resulting speedups as provisional estimates ($\dagger$).  The
estimate uses the one-pair acceptance control (293.88~s end-to-end and 265.23~s
DiT time) with the same model and sampling configuration.  CLIP-T and VBench-Q
are reported on their native $[0,1]$ scale.  Timing and quality averages are over
successful generated requests; a miss is an audit record and does not trigger a
dense fallback under this protocol.

\begin{table}[htbp]
\centering
\caption{Provisional main results on Wan2.2-TI2V-5B (121 frames, $1280\times704$,
50 UniPC steps).  $\dagger$ denotes an estimate anchored to the single-pair Dense
acceptance control; all other entries are measured in the 300-request run.  The
Carnator row uses the current unbounded-gate B5 diagnostic run (all four signature
thresholds set to $10^6$), not the final calibrated gate.}
\label{tab:main_results}
\resizebox{\linewidth}{!}{
\begin{tabular}{l c c c c c c c c c c}
\toprule
Method & Completed/Errors & CLIP-T$\uparrow$ & VBench-Q$\uparrow$
& E2E (s)$\downarrow$ & Hit E2E (s)$\downarrow$
& DiT (s)$\downarrow$ & Hit DiT (s)$\downarrow$
& E2E speedup$\uparrow$ & DiT speedup$\uparrow$ & Hit rate \\
\midrule
Dense C0$^\dagger$ & 300/0 & 0.3372 & 0.6331 & 293.88 & N/A & 265.23 & N/A & 1.00$\times$ & 1.00$\times$ & N/A \\
NIRVANA-VID & 300/0 & 0.2855 & 0.6487 & 256.09 & 256.09 & 204.21 & 204.21 & 1.15$^\dagger$ & 1.30$^\dagger$ & 80.67\% \\
TeaCache & 300/0 & 0.2938 & 0.6049 & 135.38 & N/A & 118.46 & N/A & 2.17$^\dagger$ & 2.24$^\dagger$ & N/A \\
Chorus & 259/41 & 0.2993 & 0.6407 & 264.47 & 264.47 & 175.48 & 175.48 & 1.11$^\dagger$ & 1.51$^\dagger$ & 77.61\% \\
Chorus2 & 300/0 & 0.3017 & 0.6417 & 222.44 & 222.44 & 201.58 & 201.58 & 1.32$^\dagger$ & 1.32$^\dagger$ & 80.67\% \\
Carnator (B5; unbounded gate) & 300/0 & 0.2625 & 0.6397 & 204.30 & 204.30 & 159.65 & 159.65 & 1.44$^\dagger$ & 1.66$^\dagger$ & 78.33\% \\
\bottomrule
\end{tabular}
}
\end{table}

NIRVANA-VID, Chorus, and Chorus2 use the same CLIP/Chroma threshold and consequently
have similar recall.  TeaCache has the lowest measured latency because it is a
target-only step-skipping method, but its VBench-Q is also the lowest (0.6049).
Carnator is slower than TeaCache but gives the strongest measured DiT reduction
among the cross-request methods: 159.65~s per generated request, or 1.66$\times$
under the provisional dense reference.  Its VBench-Q (0.6397) is within 0.006 of
the dense acceptance score, while CLIP-T is lower; this gap is a useful warning
that the current all-gate Carnator run is not yet a final quality-certified result.
The 14B distilled and CogVideoX configurations from the original draft are not
listed because no compatible artifacts are available in the current workspace;
they should be added only after their model-specific controls have completed.
Chorus additionally incurs 56.61~s/request of external LLM and segmentation work
and has 41 terminal errors.  Its reported 77.61\% hit rate is 201/259 successful
requests, following the aggregation protocol that excludes terminal errors from
the denominator.

The independent 694-request retrieval replay provides a complementary view of the
generation-level gate (Table~\ref{tab:retrieval_replay}).  Carnator's strict gate
accepts far fewer pairs than text-only retrieval, which is expected: it trades
recall for a lower risk of reusing incompatible object geometry or appearance.

\begin{table}[htbp]
\centering
\caption{Cross-request retrieval replay on 694 VBench prompts.  These are retrieval-only
diagnostics with an assumed step-10 cache and are not end-to-end generation results.}
\label{tab:retrieval_replay}
\begin{tabular}{l r r r r r}
\toprule
Method & Requests & Hits & Hit rate & Retrieval (ms/request) & Final cache \\
\midrule
NIRVANA-VID & 694 & 540 & 77.81\% & 12.01 & 154 \\
Chorus & 694 & 540 & 77.81\% & 12.61 & 154 \\
Carnator & 694 & 79 & 11.38\% & 70.47 & 615 \\
CHAI (reference) & 694 & 587 & 84.58\% & 8.04 & 101 \\
\bottomrule
\end{tabular}
\end{table}

\subsection{Ablation Studies}

The following tables are planning values for the unfinished ablations.  The
measured Carnator row is retained as an anchor, while entries marked $\dagger$
are idealized estimates obtained by varying one mechanism at a time around that
anchor.  They are intended to make the expected trends explicit and must be
replaced by the corresponding controlled runs before submission.

\begin{table}[htbp]
    \centering
    \caption{Expected effect of the DiT block range used for latent reuse in Carnator.
    Interval endpoints are inclusive; $\dagger$ denotes an idealized estimate.}
    \label{tab:carnator_latent_reuse_block_ablation}
    \resizebox{\linewidth}{!}{
    \begin{tabular}{c c c c c c}
        \toprule
        DiT Blocks for Latent Reuse & Speedup$\uparrow$ & E2E Time$\downarrow$
        & Denoising Time$\downarrow$ & CLIP-T$\uparrow$ & VBench-Q$\uparrow$ \\
        \midrule
        All (B5, measured) & 1.44 & 204.30 & 159.65 & 0.2625 & 0.6397 \\
        $0$--$29$ & 1.44$^\dagger$ & 204.5$^\dagger$ & 159.9$^\dagger$ & 0.263$^\dagger$ & 0.640$^\dagger$ \\
        $0$--$20$ & 1.39$^\dagger$ & 211.0$^\dagger$ & 166.0$^\dagger$ & 0.267$^\dagger$ & 0.639$^\dagger$ \\
        $5$--$25$ & 1.37$^\dagger$ & 215.0$^\dagger$ & 170.0$^\dagger$ & 0.270$^\dagger$ & 0.638$^\dagger$ \\
        $10$--$29$ & 1.41$^\dagger$ & 208.5$^\dagger$ & 163.5$^\dagger$ & 0.266$^\dagger$ & 0.639$^\dagger$ \\
        \bottomrule
    \end{tabular}
    }
\end{table}

\begin{table}[htbp]
    \centering
    \caption{Expected effect of the sparse-attention mass coverage threshold in Carnator.
    The B5 row is measured; other rows are idealized estimates ($\dagger$).}
    \label{tab:carnator_sparse_attention_ablation}
    \resizebox{\linewidth}{!}{
    \begin{tabular}{c c c c c c}
        \toprule
        Mass Coverage $\tau$ & Actual Mass Coverage & Mask Density
        & Denoising Time$\downarrow$ & E2E Time$\downarrow$ & VBench-Q$\uparrow$ \\
        \midrule
        $0.95$ (B5, measured) & 0.961 & 0.332 & 159.65 & 204.30 & 0.6397 \\
        $0.90$ & 0.915$^\dagger$ & 0.287$^\dagger$ & 153.7$^\dagger$ & 198.2$^\dagger$ & 0.637$^\dagger$ \\
        $0.85$ & 0.869$^\dagger$ & 0.245$^\dagger$ & 148.9$^\dagger$ & 192.8$^\dagger$ & 0.630$^\dagger$ \\
        \bottomrule
    \end{tabular}
    }
\end{table}

\begin{table}[htbp]
    \centering
    \caption{Component ablation for Carnator.  Component times are seconds per
    generated request; Full is measured and the remaining rows are idealized estimates ($\dagger$).}
    \label{tab:carnator_component_ablation}
    \resizebox{\textwidth}{!}{
    \begin{tabular}{l c c c c c c c c}
        \toprule
        Variant & Self-Attention & Cross-Attention & FFN & Denoising & E2E
        & Hit Rate & CLIP-T & VBench-Q \\
        \midrule
        Carnator (Full, measured) & 40.71 & 10.75 & 18.75 & 159.65 & 204.30 & 78.33\% & 0.2625 & 0.6397 \\
        w/o Sparse Attention Mask Reuse & 76.0$^\dagger$ & 10.8$^\dagger$ & 18.8$^\dagger$ & 194.0$^\dagger$ & 238.0$^\dagger$ & 78.33\%$^\dagger$ & 0.264$^\dagger$ & 0.641$^\dagger$ \\
        w/o Latent Reuse & 43.0$^\dagger$ & 11.0$^\dagger$ & 19.0$^\dagger$ & 190.0$^\dagger$ & 239.0$^\dagger$ & 78.33\%$^\dagger$ & 0.276$^\dagger$ & 0.638$^\dagger$ \\
        w/o W8A8 FFN & 40.7$^\dagger$ & 10.8$^\dagger$ & 25.0$^\dagger$ & 166.0$^\dagger$ & 211.0$^\dagger$ & 78.33\%$^\dagger$ & 0.263$^\dagger$ & 0.640$^\dagger$ \\
        \bottomrule
    \end{tabular}
    }
\end{table}

The predicted ablation ordering follows the mechanism costs.  Removing sparse
attention primarily increases self-attention time, removing latent reuse restores
target computation over the reused region, and removing W8A8 affects the FFN
column while leaving quality nearly unchanged.  The threshold sweep illustrates
the expected speed--coverage trade-off: lowering $\tau$ reduces mask density but
also reduces the mass retained by the historical connectivity.  The selected
Early Signature tuple from the current offline tuning pass
$(2.832,2.169,1.000,1.999)$ accepts 212/235
unbounded-gate hits (90.21\%) and has mean VBench-Q 0.6398, but its quality gate
has only 22 joint passes.  We therefore treat it as a diagnostic prediction and
will report the final certified tuple only after the independent calibration and
Dense C0 control complete.
